\subsection{From the lattice Lax equation to the continuum zero-curvature equation}
\label{sec:lattice-to-continuum-zc}

We now apply the lower-symbol operation to the finite-lattice equation
\eqref{eq:discrete-zero-curvature}.  At this stage the lower symbol is taken at
a fixed point of the coherent-state phase space.  Applying it to the operator
identity gives the exact relation
\begin{align}
\bigl(\partial_{t_{\rm lat}}L_{a,n}(u)\bigr)^{\downarrow}
={}&
\bigl(A_{a,n+1}(u)L_{a,n}(u)\bigr)^{\downarrow}
\nonumber\\
&-
\bigl(L_{a,n}(u)A_{a,n}(u)\bigr)^{\downarrow}.
\label{eq:lower-symbol-discrete-zc}
\end{align}
For each model, we derive the continuum Hamiltonian flow independently and
use it to identify the continuum limit of the left-hand side with the time
derivative of the classical spatial Lax matrix.  The factors in each
product on the right-hand side act on a common physical site.  The lower
symbol of each product contains a contribution in addition to the product
of the separate lower symbols.  We denote these contributions by the
auxiliary-space matrices
\begin{align}
\mathcal D^{(+)}_{a,n}(u)
:={}&
\bigl(A_{a,n+1}(u)L_{a,n}(u)\bigr)^{\downarrow}
-
A_{a,n+1}^{\downarrow}(u)L_{a,n}^{\downarrow}(u),
\label{eq:Rplus-definition}\\
\mathcal D^{(-)}_{a,n}(u)
:={}&
\bigl(L_{a,n}(u)A_{a,n}(u)\bigr)^{\downarrow}
-
L_{a,n}^{\downarrow}(u)A_{a,n}^{\downarrow}(u).
\label{eq:Rminus-definition}
\end{align}
These matrices measure the failure of the lower symbol to preserve the
operator product at the overlapping physical site.  This nonmultiplicativity
is standard in coherent-state symbol calculus and is commonly encoded by a
Berezin star product \cite{Schlichenmaier2010}.

With the definitions \eqref{eq:Rplus-definition} and
\eqref{eq:Rminus-definition}, the lower-symbol form of the lattice Lax
equation is the exact identity
\begin{align}
\bigl(\partial_{t_{\rm lat}}L_{a,n}(u)\bigr)^{\downarrow}
={}&
A_{a,n+1}^{\downarrow}(u)L_{a,n}^{\downarrow}(u)
-
L_{a,n}^{\downarrow}(u)A_{a,n}^{\downarrow}(u)
\nonumber\\
&+
\mathcal D^{(+)}_{a,n}(u)-\mathcal D^{(-)}_{a,n}(u).
\label{eq:exact-lower-symbol-lattice-zc}
\end{align}
For operators with disjoint physical support, the product coherent state
\eqref{eq:product-coherent-state} makes the corresponding difference vanish.
The difference terms in \eqref{eq:exact-lower-symbol-lattice-zc} can survive because
of the overlap described at the end of Section~\ref{sec:quantum-lattice-setup}.

We next state the common continuum scaling used by the two models studied in
this paper.  The model-dependent spectral scaling and scalar normalization of
the spatial Lax operator have already been represented abstractly by
\eqref{eq:normalized-lattice-Lax}.  We also use the freedom described below
\eqref{eq:discrete-zero-curvature} to shift the finite-lattice time Lax
operator by a scalar function times the identity.  We denote the resulting operator by
$\widehat{A}_{a,n}$.

For both models, the continuum time coordinate $t$ is related to the quantum
lattice time $t_{\rm lat}$ by
\begin{equation}
t=\epsilon^2 t_{\rm lat}.
\label{eq:common-time-scaling}
\end{equation}
Equation~\eqref{eq:common-time-scaling} fixes the normalization of the
continuum time used in this paper.  Since the scalar normalization of
$\widehat L_{a,n}$ and the scalar shift defining
$\widehat{A}_{a,n}$ are independent of $t_{\rm lat}$ and of the
lattice site, the normalized operators obey the same discrete zero-curvature
equation as $L_{a,n}$ and $A_{a,n}$.
After the model-dependent spectral scaling has been imposed, the lower symbols
of the normalized spatial and time Lax operators have the form
\begin{align}
\widehat L_{a,n}^{\downarrow}
&=
\Id_{\mathcal V_a}
+\epsilon U(x_n,t;\lambda)
+\mathcal O(\epsilon^2),
\label{eq:L-lower-common-scaling}\\
\widehat{A}_{a,n}^{\downarrow}
&=
\epsilon^2 V^{\downarrow}(x_n,t;\lambda)
+\mathcal O(\epsilon^3).
\label{eq:A-lower-common-scaling}
\end{align}
Here $V^{\downarrow}(x,t;\lambda)$ denotes the continuum matrix obtained
directly from the lower symbol of the finite-lattice time Lax operator.  The scalings \eqref{eq:common-time-scaling}--\eqref{eq:A-lower-common-scaling}
are checked for each explicit model.

The part of \eqref{eq:exact-lower-symbol-lattice-zc} containing only products
of lower symbols now has a standard continuum limit.  Using
$x_{n+1}=x_n+\epsilon$ together with
\eqref{eq:L-lower-common-scaling} and
\eqref{eq:A-lower-common-scaling}, one finds
\begin{align}
&\widehat{A}_{a,n+1}^{\downarrow}
 \widehat L_{a,n}^{\downarrow}
-
 \widehat L_{a,n}^{\downarrow}
 \widehat{A}_{a,n}^{\downarrow}
\nonumber\\
&\qquad=
\epsilon^3
\left(
\partial_x V^{\downarrow}
+[V^{\downarrow},U]
\right)_{(x_n,t;\lambda)}
+\mathcal O(\epsilon^4).
\label{eq:factorized-continuum-zc}
\end{align}
The part containing products of lower symbols starts at order
$\epsilon^3$.  The coefficient of the lower symbol of the Heisenberg
derivative at the same order is determined from the continuum Hamiltonian and
its equations of motion for each model.  This dynamical step is separate from
the algebraic treatment of the operator products.

For the normalized operators, define the corresponding differences by
\begin{align}
\widehat{\mathcal D}^{(+)}_{a,n}
:={}&
\bigl(\widehat{A}_{a,n+1}\widehat L_{a,n}\bigr)^{\downarrow}
-
\widehat{A}_{a,n+1}^{\downarrow}
\widehat L_{a,n}^{\downarrow},
\label{eq:Rplus-hat-definition}\\
\widehat{\mathcal D}^{(-)}_{a,n}
:={}&
\bigl(\widehat L_{a,n}\widehat{A}_{a,n}\bigr)^{\downarrow}
-
\widehat L_{a,n}^{\downarrow}
\widehat{A}_{a,n}^{\downarrow}.
\label{eq:Rminus-hat-definition}
\end{align}
Their continuum contribution, when the limit exists, is
\begin{equation}
\mathcal C(x,t;\lambda)
:=
\lim_{\epsilon\to0}
\frac{
\widehat{\mathcal D}^{(+)}_{a,n}
-
\widehat{\mathcal D}^{(-)}_{a,n}
}{\epsilon^3},
\qquad x=x_n.
\label{eq:C-continuum-definition}
\end{equation}
The evaluation of this limit requires both the operator expansion and
the Taylor expansion of the spin arguments.  Let $m\in\{5,6\}$ label
the model, and write $\mathcal C_m$ for its continuum contribution.
For independent spin vectors $\boldsymbol S_L$ and
$\boldsymbol S_R$ on the left and right physical sites, define the
auxiliary-space matrix functions
$\mathcal D^{(\pm)}_{m,2}(\boldsymbol S_L,\boldsymbol S_R)$ and
$\mathcal D^{(\pm)}_{m,3}(\boldsymbol S_L,\boldsymbol S_R)$ as the coefficients
\begin{align}
\widehat{\mathcal D}^{(+)}_{m;a,n}
&=\epsilon^2\mathcal D^{(+)}_{m,2}
  (\boldsymbol S_n,\boldsymbol S_{n+1})
 +\epsilon^3\mathcal D^{(+)}_{m,3}
  (\boldsymbol S_n,\boldsymbol S_{n+1})
 +\mathcal O(\epsilon^4),
\label{eq:overlap-plus-coefficients}\\
\widehat{\mathcal D}^{(-)}_{m;a,n}
&=\epsilon^2\mathcal D^{(-)}_{m,2}
  (\boldsymbol S_{n-1},\boldsymbol S_n)
 +\epsilon^3\mathcal D^{(-)}_{m,3}
  (\boldsymbol S_{n-1},\boldsymbol S_n)
 +\mathcal O(\epsilon^4).
\label{eq:overlap-minus-coefficients}
\end{align}
The dependence on the spectral parameter and continuum coupling is
suppressed.  The final subscript denotes the power of $\epsilon$ in
the operator-product difference before expanding the spin arguments.
For both models, the quadratic coefficients agree at coincident spins,
$\mathcal D^{(+)}_{m,2}(\boldsymbol S,\boldsymbol S)
=\mathcal D^{(-)}_{m,2}(\boldsymbol S,\boldsymbol S)$,
as verified in the model calculations below.  Taylor expansion at
$x=x_n$ then gives
\begin{align}
\mathcal C_m
={}&\partial_x S^i
\left.
\left(
\frac{\partial\mathcal D^{(+)}_{m,2}}{\partial S_R^i}
+\frac{\partial\mathcal D^{(-)}_{m,2}}{\partial S_L^i}
\right)\right|_{\boldsymbol S_L=\boldsymbol S_R=\boldsymbol S}
\nonumber\\
&+\mathcal D^{(+)}_{m,3}(\boldsymbol S,\boldsymbol S)
 -\mathcal D^{(-)}_{m,3}(\boldsymbol S,\boldsymbol S).
\label{eq:overlap-source-taylor}
\end{align}
The Cartesian derivatives act on the polynomial lower symbols before
imposing the unit-spin constraints and before setting the two arguments
equal.  Both derivatives have a plus sign:
the shift $\boldsymbol S_{n-1}=\boldsymbol S-\epsilon\partial_x
\boldsymbol S+\cdots$ occurs in the term subtracted in
Eq.~\eqref{eq:C-continuum-definition}.  Equation~\eqref{eq:overlap-source-taylor}
separates the contribution from the spatial variation of the quadratic
coefficients from the cubic coefficients at coincident spins.

Once the model-specific Hamiltonian check identifies the continuum limit of
the lower symbol of the Heisenberg derivative on the left-hand side of
\eqref{eq:exact-lower-symbol-lattice-zc} with $\partial_tU$, the
order-$\epsilon^3$ equation becomes
\begin{equation}
\partial_t U
=
\partial_xV^{\downarrow}
+[V^{\downarrow},U]
+\mathcal C.
\label{eq:continuum-zc-with-C}
\end{equation}

The continuum zero-curvature equation is recovered if the matrix
$\mathcal C$ can be written in the local form
\begin{equation}
\mathcal C
=
\partial_x\Delta V+[\Delta V,U]
\label{eq:C-absorbable-condition}
\end{equation}
for some auxiliary-space matrix $\Delta V(x,t;\lambda)$.  In that case the
continuum time Lax matrix
\begin{equation}
V:=V^{\downarrow}+\Delta V
\label{eq:corrected-continuum-V-definition}
\end{equation}
satisfies
\begin{equation}
\partial_tU-\partial_xV+[U,V]=0.
\label{eq:continuum-zero-curvature}
\end{equation}
Equations \eqref{eq:C-continuum-definition}--\eqref{eq:continuum-zero-curvature}
formulate the continuum problem.  The matrix $\Delta V$ is determined from
the finite-lattice operator products for each model.
